\documentclass{article}
\textheight 23.5cm \textwidth 15.8cm
%\leftskip -1cm
\topmargin -1.5cm \oddsidemargin 0.3cm \evensidemargin -0.3cm
%\documentclass[final]{siamltex}

\usepackage{ctex}
\usepackage{verbatim}
\usepackage{fancyhdr}
\usepackage{graphicx}
\usepackage{booktabs}

%\pagestyle{fancy} \lhead{FDM Homework Template} \chead{}
%\rhead{\bfseries Yan XU}
%
%\lfoot{} \cfoot{} \rfoot{\thepage}
%\renewcommand{\headrulewidth}{0.4pt}
%\renewcommand{\footrulewidth}{0.4pt}
%
\title{第二周程序作业}
\author{孙天阳 SA23001051}
\date{2024年9月25日}
\begin{document}
\maketitle

\section{引言}
在本次作业中, 我们采用单元循环而不是基循环来组装刚度矩阵$ A $, 这种方式更易迁移到高维。
\section{方法}

\section{结果}
\begin{table}[h!]
    \centering
    \caption{\( L^2 \) 和 \( H^1 \) 误差及收敛阶数}
    \label{tab:results}
    \begin{tabular}{ccccc}
        \toprule
        \( N \) & \( L^2 \) 误差 & \( L^2 \) 收敛阶数 & \( H^1 \) 误差 & \( H^1 \) 收敛阶数 \\
        \midrule
        10  & \( 1.40854 \times 10^{-3} \) & -    & \( 4.90167 \times 10^{-2} \) & - \\
        20  & \( 3.86443 \times 10^{-4} \) & 1.87 & \( 2.56657 \times 10^{-2} \) & 0.93 \\
        40  & \( 1.01409 \times 10^{-4} \) & 1.93 & \( 1.31444 \times 10^{-2} \) & 0.97 \\
        80  & \( 2.42843 \times 10^{-5} \) & 2.06 & \( 6.65316 \times 10^{-3} \) & 0.98 \\
        \bottomrule
    \end{tabular}
\end{table}
\section{讨论}
可以看到结果与第一周程序作业完全一样，这是符合预期的，因为实际上每个量都没有发生改变，改变的只是得到这些量的方式。
\end{document}
